\documentclass[border=3pt]{standalone}
\usepackage[european]{circuitikz}
\usetikzlibrary{calc}
\standaloneenv{circuitikz}

\begin{document}
\begin{circuitikz}[
  scale=1,
  european,
  wire/.style={line width=0.7pt},
  transistor/.style={line width=0.7pt},
  rail/.style={line width=1.4pt},
  netlabel/.style={font=\ttfamily\footnotesize},
  devlabel/.style={font=\ttfamily\footnotesize, align=left},
]

% 28 nm PMOS-input preamp, reconstructed from the FSiC 2026 reference slide.
% MP4/MP3 form the PMOS current mirror.
% MP1/MP2 are the PMOS input pair.
% MN1/MN2 are static NMOS loads biased by vbias.

\coordinate (XP)    at (0,1.45);
\coordinate (XN)    at (4,1.45);
\coordinate (PSRC)  at (2,2.65);
\coordinate (IBIAS) at (0,4.05);
\coordinate (PMIR)  at (2,4.05);

% Top PMOS current mirror. Gates point inward.
\node[pmos, transistor, xscale=-1, anchor=D] (mp4) at (IBIAS) {};
\node[pmos, transistor, anchor=D] (mp3) at (PMIR) {};

% PMOS input pair. Gates point outward to inn/inp.
% xp/xn are centered between the PMOS and NMOS drain connection points.
\node[pmos, transistor, anchor=D] (mp1) at ($(XP)+(0,0.28)$) {};
\node[pmos, transistor, xscale=-1, anchor=D] (mp2) at ($(XN)+(0,0.28)$) {};

% NMOS loads. Gates point inward and share vbias.
\node[nmos, transistor, xscale=-1, anchor=D] (mn1) at ($(XP)+(0,-0.28)$) {};
\node[nmos, transistor, anchor=D] (mn2) at ($(XN)+(0,-0.28)$) {};

% VDD rails for current mirror.
\draw[wire]
  (mp4.S) -- ++(0,0.35) coordinate (vdd_mp4)
  (mp3.S) -- ++(0,0.35) coordinate (vdd_mp3);
\draw[rail]
  ($(vdd_mp4)+(-0.45,0)$) -- ($(vdd_mp4)+(0.45,0)$)
  ($(vdd_mp3)+(-0.45,0)$) -- ($(vdd_mp3)+(0.45,0)$);

% MP4 diode connection and PMOS mirror gate connection. Orthogonal only.
\coordinate (IBNODE) at ($(mp4.D)+(0,-0.35)$);
\draw[wire]
  (mp4.D) -- (IBNODE)
  (IBNODE) -| (mp4.G)
  (mp4.G) -- (mp3.G);

% Bias input enters from the left of the diode-connected device.
\draw[wire]
  (IBNODE) -- ++(-0.60,0) node[netlabel,left]{ibias};

% Mirror output biases the shared PMOS input-pair source.
% Keep the input-pair source bus straight and no longer than the two source anchors.
\draw[wire]
  (mp1.S) -- (mp2.S)
  (mp3.D) -- ($(mp1.S)!0.5!(mp2.S)$);

% Output nodes and short stubs.
\draw[wire]
  (XP) node[circ]{} -- ++(-0.22,0) node[netlabel,left]{xp}
  (XN) node[circ]{} -- ++(0.22,0) node[netlabel,right]{xn};

% PMOS input labels at the gate anchors.
\draw[wire]
  (mp1.G) -- ++(-0.35,0) node[netlabel,left]{inn}
  (mp2.G) -- ++(0.35,0) node[netlabel,right]{inp};

% Vertical output/load paths.
\draw[wire]
  (mp1.D) -- (mn1.D)
  (mp2.D) -- (mn2.D);

% NMOS load sources to ground.
\draw[wire]
  (mn1.S) -- ++(0,-0.20) node[ground]{}
  (mn2.S) -- ++(0,-0.20) node[ground]{};

% NMOS load gates share vbias on one centered horizontal bus.
\draw[wire]
  (mn1.G) -- (mn2.G) node[netlabel,midway,below]{vbias};

% Device annotations. Keep clumps left-aligned and outside signal routes.
\node[devlabel, anchor=east] at ($(mp4)+(-0.32,0.20)$) {MP4\\w=120n\\l=100n};
\node[devlabel, anchor=west] at ($(mp3)+(0.32,0.20)$) {MP3\\w=120n\\l=100n};

\node[devlabel, anchor=west] at ($(mp1)+(0.35,-0.16)$) {MP1\\w=240n\\l=50n};
\node[devlabel, anchor=east] at ($(mp2)+(-0.35,-0.16)$) {MP2\\w=240n\\l=50n};

\node[devlabel, anchor=east] at ($(mn1)+(-0.32,-0.08)$) {MN1\\w=60n\\l=50n};
\node[devlabel, anchor=west] at ($(mn2)+(0.32,-0.08)$) {MN2\\w=60n\\l=50n};

\end{circuitikz}
\end{document}
